\scanpage{40}
% archive sheet number: 39

\uline{8)} Anneau \unsure{irreductible|irréductible}

\[
 A_{\rm red}\overset{\rm df}{=}A/\underbrace{\varinjlim_n (I_n}_{\text{nilradical }N}=\operatorname{Ker}(A\xrightarrow{x\mapsto x^n}A))
 \quad\text{est intègre}
\]

i.e.\quad a)
\diagram{\begin{tikzcd}[ampersand replacement=\&]
 A\times N\amalg N\times A\arrow[r,"v"]\& D_N\arrow[r,hook]\arrow[d]\& A\times A\arrow[d,"{(x,y)\mapsto xy}"]\\
 \text{$v$ épi}\&N\arrow[r,hook]\arrow[ur,phantom,"\rm cart"]\& A
\end{tikzcd}}

\[
\begin{aligned}
 \mathrm b)\qquad&V(1_A)\simeq\varnothing_{\mathcal E}\\
 \deleted{\unsure{\mathrm{i.e.}\ \forall f}}\quad\text{i.e.}\quad&
 \left\{\begin{aligned}
 &\forall f,g\in A(U),\quad\varinjlim_n V(f^ng^n)\simeq\varinjlim V(f^n)\vee\varinjlim V(g^n)\\
 &V(1_A)\simeq\varnothing_{\mathcal E}
 \end{aligned}\right.
\end{aligned}
\]

$X_{\rm irr}\simeq\widetilde R_{\mathrm{top\,irr}}$\quad top irred la topologie où une famille
\deleted{$X_i\to$}$X$ est couvrante \uline{ssi} $\forall$ comp. \deleted{sont} irréductible \deleted{$Y$} de $X$,
et \deleted{\unsure{la}}\unsure{ss-schéma} \unsure{induit} $T_i$, l'un des $X_i$ a une section sur $T_{\unsure{i\,\rm red|i|\rm red}}$.

On obtient un exemple d'un sous-topos $X_{\rm irr}$ non \deleted{\unsure{quasi-compact}}\inserted{cohérent} du topos cohérent $X_{\rm ann}$.

\uline{NB} Cette topologie n'est pas quasi-compacte, il y a des recouvrements
(de \deleted{\unsure{schémas}} tout schéma \deleted{\unsure{dont les}}\inserted{ayant une} comp. connexe non \deleted{\unsure{quasi}} irréductible)
qui n'admettent pas de ss-rec$^t$ fini

\uline{NB} Les $F\in\operatorname{Ob}X_{\rm irr}\subset\operatorname{Ob}\widehat R$ transforment sommes finies en produits
\inserted{$\operatorname{Ob}\widetilde R_{\mathrm{top\,irr}}$}

\uline{9)} \uline{Anneau connexe} (i.e. à sp. conn.)\quad $X_{\rm conn}=\widetilde R_{\mathrm{top\,conn}}\simeq\widehat R_{\rm conn}$
($R_{\rm conn}$ ss-cat des sch. aff. conn.)

$\forall f\in A(U)$, $f^2=f\Rightarrow f=1$ ou $f=0$ loc.
Suffit dans cas universel $U=A$;

top. où les familles couvrantes $X_i\to X$ sont celles qui ont une section sur les composantes connexes \inserted{\unsure{c.à.d. ens. d'inclusions|cat. ens. d'inclusions} \unsure{crible|vide} si $X=\varnothing$}.

\diagram{\begin{tikzcd}[ampersand replacement=\&,column sep=small,row sep=small]
 \& A\arrow[r,"{x\mapsto x^2-x}"]\arrow[d]\& A\\
 e\amalg e\arrow[r,"u"]\&\Pi\arrow[r]\& e\arrow[u,"0_A"']\&\text{$u$ épi}\\
 \deleted{\unsure{0,1}}\&\text{(projecteurs)}
\end{tikzcd}}

$\widetilde R_{\rm conn}$ \unsure{c'est} préfaisceaux $R^\circ\to(\mathrm{Ens})$ transformant sommes
\deleted{\unsure{(ensemblistes) en produits}} de deux (donc $n\geq 1$) objets en produits
(\uline{NB} si \unsure{seulement} connexes \deleted{\unsure{conn}} mais \unsure{doit} $\ne\varnothing$
on trouverait les foncteurs \unsure{$F$} qui transforment sommes finies en \unsure{sommes}, \unsure{mais}
\unsure{n'envoyant pas|n'envoyons pas} nécess$^t$ \unsure{d'ensemble|d'ens. vide}. O.K.).

\[
 X_{\mathrm{conn}\ne\varnothing}\simeq\widehat R_{\mathrm{conn}\ne\varnothing}
\]

\uline{10)} Anneau \deleted{\unsure{compact}}\inserted{local de dim 0} \deleted{\unsure{(i.e. à spectre compact)}}

\[
\begin{aligned}
 &\forall f\in A(U),\qquad U_f\amalg_{\unsure{R|B}}\varinjlim V(f^n)\longrightarrow U\quad\text{épi}\\
 &\text{Suffit ds cas universel}\qquad A^*\amalg_{\unsure{R|A}}N_A\deleted{\unsure{\rm loc}}
 \longrightarrow A\quad\text{est épi}\quad
 (N_A:=\varinjlim_n(\operatorname{Ker}(A\xrightarrow{x\mapsto x^n}A))
\end{aligned}
\]
